% Read-only excerpt from the current adjacent chapter source V5-C03.tex.
% Source lines 200--232; only the portion needed to adjudicate FIG-P600-01 is retained.
\paragraph{读前自检：细致平衡。}设$\pi$为概率密度，马尔可夫核$K$有密度$k$，且$\pi(x)k(x,y)=\pi(y)k(y,x)$对$\mu\otimes\mu$几乎处处成立。请对$x$积分，核对总流入为$\pi(y)$，故$\pi K=\pi$。\par

\begin{definition}[细致平衡]
若对任意可测集合$A,B\in\mathcal B$都有
\begin{equation}\label{eq:V5-C03-general-detailed-balance}
\int_A\pi(\mathrm dx)K(x,B)
=\int_B\pi(\mathrm dy)K(y,A),
\end{equation}
则称$K$关于$\pi$满足细致平衡，或称平稳链关于$\pi$可逆。若$K$有密度$k(x,y)$，条件化为
\begin{equation}\label{eq:V5-C03-density-detailed-balance}
\pi(x)k(x,y)=\pi(y)k(y,x)
\quad\text{对}\ \mu\otimes\mu\ \text{几乎处处成立}.
\end{equation}
\end{definition}

从局部流到平稳分布的推导分三步。第一步，式\eqref{eq:V5-C03-density-detailed-balance}把$x\to y$的平稳概率流与$y\to x$的反向流配对。第二步，对$x$积分，
\[
\int\pi(x)k(x,y)\,\mu(\mathrm dx)
=\int\pi(y)k(y,x)\,\mu(\mathrm dx).
\]
第三步使用$\int k(y,x)\mu(\mathrm dx)=1$，右端成为$\pi(y)$，于是总流入等于目标质量。图\ref{fig:V5-C03-mh-balance-flux}把成对流与拒绝自环分开绘制。
\input{../../绘图源码/第05册_采样方法主题模型与图排序/V5-C03/fig_v5_c03_mh_balance_flux.tex}

\begin{proposition}[细致平衡推出平稳]\label{prop:V5-C03-db-stationary}
若概率测度$\pi$与马尔可夫核$K$满足式\eqref{eq:V5-C03-general-detailed-balance}，则$\pi K=\pi$。
\end{proposition}
